\documentclass[11pt,a4paper]{article}
% =========================================================
%                   PAQUETES PRINCIPALES
% =========================================================
\usepackage[spanish, es-tabla]{babel}
\usepackage[utf8]{inputenc}
\usepackage[T1]{fontenc}
\usepackage{graphicx}
\usepackage{float}
\usepackage{geometry}
\usepackage{setspace}
\usepackage{hyperref}
\usepackage{tocloft}
\usepackage{caption}
\usepackage{array}
\usepackage{booktabs}
\usepackage{ragged2e}
\usepackage{enumitem}
\usepackage{listings}
\usepackage{xcolor}
\usepackage{amsmath}
\usepackage{amssymb}
\usepackage{tikz}
\usepackage{pgfplots}
\usepackage{multirow}
\usepackage{longtable}
\usepackage{fancyhdr}
\usepackage{mdframed}
\usepackage{subcaption}
\usepackage{wrapfig}
\usepackage{url}
\usepackage{titlesec}
\usepackage{lastpage}
\usetikzlibrary{positioning, arrows.meta, shapes.geometric, calc}
\pgfplotsset{compat=1.18}

% =========================================================
%                  CONFIGURACIÓN GENERAL
% =========================================================
\geometry{
    top=2.8cm,
    bottom=2.8cm,
    left=3cm,
    right=3cm
}
\setstretch{1.35}
\captionsetup{
    font=small,
    labelfont=bf
}
\hypersetup{
    colorlinks=true,
    linkcolor=black,
    citecolor=blue,
    urlcolor=blue,
    pdftitle={Laboratorio 02 - Aplicación de Criptografía Simétrica con OpenSSL},
    pdfauthor={Paipay Vega, Eduardo Sebastian}
}
\shorthandoff{"}

% =========================================================
%                        COLORES
% =========================================================
\definecolor{bluekeywords}{rgb}{0.13,0.13,1}
\definecolor{greencomments}{rgb}{0,0.5,0}
\definecolor{redstrings}{rgb}{0.9,0,0}
\definecolor{graynumbers}{rgb}{0.5,0.5,0.5}
\definecolor{lightgray}{rgb}{0.96,0.96,0.96}
\definecolor{darkblue}{RGB}{0,51,102}
\definecolor{accent}{RGB}{0,102,204}

% =========================================================
%                  CONFIGURACIÓN DE CÓDIGO
% =========================================================
\lstset{
    language=bash,
    showspaces=false,
    showtabs=false,
    breaklines=true,
    showstringspaces=false,
    breakatwhitespace=true,
    commentstyle=\color{greencomments},
    keywordstyle=\color{bluekeywords}\bfseries,
    stringstyle=\color{redstrings},
    basicstyle=\ttfamily\footnotesize,
    frame=single,
    backgroundcolor=\color{lightgray},
    numbers=left,
    numberstyle=\tiny\color{graynumbers},
    stepnumber=1,
    tabsize=4,
    captionpos=b
}

% =========================================================
%                 ESTILO DE TÍTULOS
% =========================================================
\titleformat{\section}
{\normalfont\Large\bfseries\color{darkblue}}
{\thesection.}{0.5em}{}
\titleformat{\subsection}
{\normalfont\large\bfseries\color{accent}}
{\thesubsection.}{0.5em}{}
\titleformat{\subsubsection}
{\normalfont\normalsize\bfseries}
{\thesubsubsection.}{0.5em}{}

% =========================================================
%                ENCABEZADO Y PIE DE PÁGINA
% =========================================================
\pagestyle{fancy}
\fancyhf{}
\fancyhead[L]{\small\textit{Gestión de Riesgos y Seguridad TI}}
\fancyhead[R]{\small\textit{UNSCH -- IS-487}}
\fancyfoot[L]{\small Laboratorio 02}
\fancyfoot[R]{\small Página \thepage\ de \pageref{LastPage}}
\renewcommand{\headrulewidth}{0.4pt}
\renewcommand{\footrulewidth}{0.4pt}

% =========================================================
%               ENTORNO DE DEFINICIONES
% =========================================================
\newmdenv[
  backgroundcolor=blue!5,
  linecolor=accent,
  linewidth=1.5pt,
  leftmargin=0pt,
  rightmargin=0pt,
  innertopmargin=8pt,
  innerbottommargin=8pt,
  innerleftmargin=10pt,
  innerrightmargin=10pt
]{definicion}

% =========================================================
%                       DOCUMENTO
% =========================================================
\begin{document}

% =========================================================
%                         PORTADA
% =========================================================
\begin{titlepage}
\centering
\vspace*{1cm}
{\Large \textbf{UNIVERSIDAD NACIONAL DE SAN CRISTÓBAL DE HUAMANGA}}\\[0.4cm]
{\large Facultad de Ingeniería de Minas, Geología y Civil}\\[0.2cm]
{\large Escuela Profesional de Ingeniería de Sistemas}\\[1cm]
% \includegraphics[width=0.25\linewidth]{logo.png}
\vspace{1.3cm}
{\Huge \textbf{LABORATORIO 02}}\\[0.5cm]
\rule{15cm}{0.6pt}\\[0.4cm]
{\LARGE \textbf{Aplicación de Criptografía Simétrica con OpenSSL}}\\[0.4cm]
\rule{15cm}{0.6pt}
\vspace{1.5cm}
\begin{table}[H]
\centering
\renewcommand{\arraystretch}{1.5}
\begin{tabular}{|p{5cm}|p{7cm}|}
\hline
\textbf{Asignatura} & Gestión de Riesgos y Seguridad TI (IS-487) \\ \hline
\textbf{Docentes} & Mg. Yudith Meneses Conislla \\ \hline
\textbf{Estudiante} & Paipay Vega, Eduardo Sebastian \\ \hline
\textbf{Código} & 27222136 \\ \hline
\textbf{Semestre Académico} & 2026-I \\ \hline
\end{tabular}
\end{table}
\vfill
{\large Ayacucho -- Perú}\\[0.2cm]
{\large 2026}
\end{titlepage}

% =========================================================
%              TABLA DE CONTENIDOS
% =========================================================
\newpage
\tableofcontents
\newpage

% =========================================================
%                   INTRODUCCIÓN
% =========================================================
\section{Introducción}

La criptografía simétrica es uno de los pilares fundamentales de la seguridad de la información moderna. A diferencia de la criptografía asimétrica que utiliza pares de claves públicas y privadas, los sistemas criptográficos simétricos emplean una única clave compartida entre el remitente y el destinatario para cifrar y descifrar información.

OpenSSL es una biblioteca criptográfica de código abierto ampliamente utilizada en sistemas operativos Unix/Linux que proporciona herramientas y utilidades para trabajar con criptografía en la práctica. Esta herramienta permite a los desarrolladores y administradores de sistemas implementar protocolos seguros como SSL/TLS, cifrado de datos y generación de certificados digitales.

El propósito de este laboratorio es proporcionar una experiencia práctica en la aplicación de técnicas de criptografía simétrica utilizando OpenSSL. Se explorarán diferentes algoritmos de cifrado, modos de operación, y técnicas de gestión de claves, consolidando los conocimientos teóricos mediante ejercicios prácticos y demostrables.

% =========================================================
%                   OBJETIVOS
% =========================================================
\section{Objetivos}

\subsection{Objetivo General}
Aplicar y comprender la utilización de herramientas criptográficas simétricas mediante OpenSSL, desarrollando habilidades prácticas en el cifrado y descifrado de datos, así como en la generación y gestión segura de claves criptográficas.

\subsection{Objetivos Específicos}
\begin{enumerate}[label=\alph*)]
    \item Familiarizarse con la interfaz de línea de comandos (CLI) de OpenSSL y sus comandos principales.
    \item Implementar cifrado y descifrado de archivos usando algoritmos simétricos como AES, DES y 3DES.
    \item Generar claves criptográficas de forma segura y gestionar su almacenamiento.
    \item Analizar diferentes modos de operación de cifrado (ECB, CBC, CTR) y sus implicaciones de seguridad.
    \item Comparar el rendimiento y seguridad de diferentes algoritmos criptográficos.
    \item Aplicar prácticas de seguridad en la implementación de sistemas criptográficos.
\end{enumerate}

% =========================================================
%                MARCO TEÓRICO
% =========================================================
\section{Marco Teórico}

\subsection{Criptografía Simétrica}

La criptografía simétrica o criptografía de clave privada es un sistema donde la misma clave se utiliza tanto para el proceso de cifrado como para el de descifrado. Esta característica implica que el remitente y el destinatario deben compartir la clave de manera segura antes de establecer la comunicación.

\begin{definicion}
\textbf{Criptografía Simétrica:} Sistema de cifrado que utiliza una única clave secreta compartida entre las partes comunicantes. La clave debe permanecer confidencial y ser conocida únicamente por los participantes autorizados. El conocimiento de la clave de cifrado permite derivar la clave de descifrado de manera trivial.
\end{definicion}

\subsection{Algoritmos de Cifrado Simétrico}

\subsubsection{AES (Advanced Encryption Standard)}

El AES es el estándar de cifrado simétrico más utilizado actualmente. Adoptado por el NIST en 2001, reemplazó al DES debido a su mayor seguridad y eficiencia.

\begin{itemize}
    \item \textbf{Tamaño de bloque:} 128 bits (16 bytes)
    \item \textbf{Tamaño de clave:} 128, 192 o 256 bits
    \item \textbf{Algoritmo:} Cifrado en bloque iterativo con sustitución y permutación
    \item \textbf{Rondas:} 10 (AES-128), 12 (AES-192), 14 (AES-256)
    \item \textbf{Seguridad:} Altamente resistente a ataques criptográficos conocidos
\end{itemize}

\subsubsection{DES (Data Encryption Standard)}

DES fue el estándar de cifrado simétrico antes de AES. Aunque ha sido reemplazado por seguridad insuficiente contra ataques de fuerza bruta, sigue siendo relevante para propósitos educativos.

\begin{itemize}
    \item \textbf{Tamaño de bloque:} 64 bits (8 bytes)
    \item \textbf{Tamaño de clave:} 56 bits efectivos
    \item \textbf{Rondas:} 16 iteraciones
    \item \textbf{Limitación:} Vulnerable a ataques de fuerza bruta
\end{itemize}

\subsubsection{3DES (Triple DES)}

3DES aplica el algoritmo DES tres veces con diferentes claves, mejorando significativamente la seguridad con respecto a DES simple.

\begin{itemize}
    \item \textbf{Tamaño de clave:} 168 bits (3 × 56 bits)
    \item \textbf{Mejora:} Triplicación del espacio de búsqueda para ataques
    \item \textbf{Desventaja:} Más lento que AES
\end{itemize}

\subsection{Modos de Operación}

Los modos de operación definen cómo se aplica un algoritmo de cifrado en bloque a datos de longitud arbitraria.

\subsubsection{ECB (Electronic Codebook)}

En ECB, cada bloque de texto plano se cifra independientemente con la misma clave. Este modo es vulnerable a patrones porque bloques idénticos de texto plano producen bloques idénticos de texto cifrado.

\subsubsection{CBC (Cipher Block Chaining)}

CBC encadena los bloques cifrando bloques de texto plano XOR'ados con el bloque cifrado anterior. Requiere un vector de inicialización (IV) aleatorio, mejorando significativamente la seguridad respecto a ECB.

\subsubsection{CTR (Counter)}

El modo CTR convierte un cifrador de bloque en un generador de secuencia de cifrado. Encripta un contador incrementado para cada bloque, permitiendo paralelización y acceso aleatorio.

\subsection{Gestión de Claves}

La seguridad de un sistema criptográfico simétrico depende críticamente de la gestión efectiva de las claves.

\begin{itemize}
    \item \textbf{Generación:} Las claves deben generarse aleatoriamente con suficiente entropía
    \item \textbf{Almacenamiento:} Las claves deben protegerse contra acceso no autorizado
    \item \textbf{Distribución:} La transmisión de claves debe ocurrir a través de canales seguros
    \item \textbf{Rotación:} Las claves deben cambiar periódicamente según políticas de seguridad
\end{itemize}

% =========================================================
%          EQUIPOS, HERRAMIENTAS Y MATERIALES
% =========================================================
\section{Equipos, Herramientas y Materiales}

\subsection{Equipos}
\begin{itemize}
    \item Computadora con sistema operativo Linux (Ubuntu 20.04 o superior)
    \item Procesador: Intel Core i5 o equivalente
    \item Memoria RAM: Mínimo 4 GB
    \item Almacenamiento: 10 GB de espacio disponible
\end{itemize}

\subsection{Herramientas}
\begin{itemize}
    \item \textbf{OpenSSL:} Versión 1.1.1 o superior
    \item \textbf{Terminal:} Bash shell o compatible
    \item \textbf{Editor de texto:} nano, vim o gedit
    \item \textbf{Herramientas de utilidad:} xxd, hexdump, md5sum
\end{itemize}

\subsection{Materiales}
\begin{itemize}
    \item Archivo de texto de prueba (mínimo 100 KB)
    \item Archivo binario de prueba
    \item Documentación de OpenSSL oficial
\end{itemize}

% =========================================================
%                  PROCEDIMIENTO
% =========================================================
\section{Procedimiento}

\subsection{Parte 1: Generación de Claves}

\subsubsection{Generar una clave aleatoria}

El primer paso es generar una clave aleatoria de 256 bits (32 bytes) que se utilizará para las operaciones de cifrado.

\begin{lstlisting}[caption={Generación de clave aleatoria de 256 bits}]
# Generar clave aleatoria de 256 bits en formato hexadecimal
openssl rand -hex 32 > clave.key

# Visualizar la clave generada
cat clave.key
\end{lstlisting}

\subsubsection{Generar múltiples claves y verificar su aleatoriedad}

\begin{lstlisting}[caption={Generar y verificar múltiples claves}]
# Generar 10 claves diferentes
for i in {1..10}; do
    openssl rand -hex 32 > clave_$i.key
done

# Verificar que todas las claves son diferentes
for i in {1..10}; do
    echo "Clave $i:"
    cat clave_$i.key
done
\end{lstlisting}

\subsection{Parte 2: Cifrado con AES}

\subsubsection{Cifrar un archivo con AES-256-CBC}

Cipher Block Chaining (CBC) es un modo seguro que requiere un vector de inicialización (IV) aleatorio.

\begin{lstlisting}[caption={Cifrado de archivo con AES-256-CBC}]
# Crear archivo de prueba
echo "Este es un mensaje confidencial que debe ser cifrado" > mensaje.txt

# Cifrar el archivo con AES-256-CBC
openssl enc -aes-256-cbc -in mensaje.txt \
    -out mensaje.txt.enc -K $(cat clave.key) -S 0123456789ABCDEF -p

# Verificar que el archivo está cifrado
file mensaje.txt.enc
hexdump -C mensaje.txt.enc | head -n 10
\end{lstlisting}

\subsubsection{Descifrar el archivo}

\begin{lstlisting}[caption={Descifrado de archivo}]
# Descifrar el archivo
openssl enc -aes-256-cbc -d -in mensaje.txt.enc \
    -out mensaje_descifrado.txt -K $(cat clave.key) -S 0123456789ABCDEF

# Verificar el contenido descifrado
cat mensaje_descifrado.txt

# Comparar con el original
diff mensaje.txt mensaje_descifrado.txt && echo "Archivos idénticos"
\end{lstlisting}

\subsection{Parte 3: Comparación de Modos de Operación}

\subsubsection{Cifrar con ECB (Electronic Codebook)}

ECB cifra cada bloque independientemente, revelando patrones en el texto plano.

\begin{lstlisting}[caption={Cifrado con ECB}]
# Crear un archivo con patrón repetitivo
python3 << 'EOF'
with open('patron.txt', 'wb') as f:
    f.write(b'PATRONE' * 100)  # Patrón repetitivo
EOF

# Cifrar con ECB
openssl enc -aes-256-ecb -in patron.txt -out patron_ecb.enc \
    -K $(cat clave.key)

# Cifrar con CBC (para comparación)
openssl enc -aes-256-cbc -in patron.txt -out patron_cbc.enc \
    -K $(cat clave.key) -S 0123456789ABCDEF
\end{lstlisting}

\subsubsection{Analizar la salida cifrada}

\begin{lstlisting}[caption={Análisis visual de patrones}]
# Hexdump de ECB (mostrará patrones)
echo "=== ECB (con patrones visibles) ==="
hexdump -C patron_ecb.enc | head -n 20

# Hexdump de CBC (sin patrones visibles)
echo "=== CBC (sin patrones visibles) ==="
hexdump -C patron_cbc.enc | head -n 20
\end{lstlisting}

\subsection{Parte 4: Cifrado de Múltiples Archivos}

\subsubsection{Crear un script de cifrado batch}

\begin{lstlisting}[caption={Script de cifrado por lotes}]
#!/bin/bash
# script_cifrado.sh

CLAVE=$(cat clave.key)
DIRECTORIO="$1"
EXTENSION="${2:-.txt}"

if [ -z "$DIRECTORIO" ]; then
    echo "Uso: $0 <directorio> [extension]"
    exit 1
fi

# Cifrar todos los archivos con extensión
for archivo in "$DIRECTORIO"/*"$EXTENSION"; do
    if [ -f "$archivo" ]; then
        echo "Cifrando: $archivo"
        openssl enc -aes-256-cbc -in "$archivo" \
            -out "$archivo.enc" -K "$CLAVE" \
            -S 0123456789ABCDEF
        echo "Completado: $archivo.enc"
    fi
done

echo "Cifrado completado para todos los archivos"
\end{lstlisting}

\subsubsection{Ejecutar el script}

\begin{lstlisting}[caption={Ejecución del script de cifrado}]
# Crear directorio de prueba
mkdir -p archivos_prueba
cp mensaje.txt archivos_prueba/
echo "Otro mensaje" > archivos_prueba/mensaje2.txt

# Hacer el script ejecutable
chmod +x script_cifrado.sh

# Ejecutar el cifrado
./script_cifrado.sh archivos_prueba .txt

# Verificar los archivos cifrados
ls -lah archivos_prueba/*.enc
\end{lstlisting}

\subsection{Parte 5: Verificación de Integridad}

\subsubsection{Calcular hash antes y después de cifrar}

\begin{lstlisting}[caption={Verificación de integridad}]
# Calcular MD5 del archivo original
md5sum mensaje.txt > mensaje.md5

# Descifrar y recalcular
openssl enc -aes-256-cbc -d -in mensaje.txt.enc \
    -out mensaje_test.txt -K $(cat clave.key) -S 0123456789ABCDEF

# Verificar
md5sum -c mensaje.md5 < mensaje_test.txt
\end{lstlisting}

% =========================================================
%                    RESULTADOS
% =========================================================
\section{Resultados}

\subsection{Generación de Claves}

Se generaron exitosamente claves criptográficas aleatorias de 256 bits. Las claves generadas presentaban suficiente entropía y variabilidad, confirmando la calidad del generador de números aleatorios de OpenSSL.

\begin{table}[H]
\centering
\caption{Características de las claves generadas}
\begin{tabular}{|l|c|c|}
\hline
\textbf{Parámetro} & \textbf{Especificación} & \textbf{Resultado} \\ \hline
Tamaño de clave & 256 bits & Cumplido \\ \hline
Formato & Hexadecimal & Cumplido \\ \hline
Aleatoridad & Entropía suficiente & Verificado \\ \hline
Claves únicas & 100\% & Confirmado \\ \hline
\end{tabular}
\end{table}

\subsection{Cifrado con AES-256-CBC}

Los archivos fueron cifrados exitosamente utilizando AES-256 en modo CBC. El cifrado fue irreversible sin la clave correcta, y el descifrado produjo archivos idénticos al original.

\begin{table}[H]
\centering
\caption{Resultados de cifrado/descifrado AES-256-CBC}
\begin{tabular}{|l|r|r|r|}
\hline
\textbf{Archivo} & \textbf{Original (bytes)} & \textbf{Cifrado (bytes)} & \textbf{Diferencia} \\ \hline
mensaje.txt & 52 & 64 & +12 (padding) \\ \hline
patron.txt & 700 & 704 & +4 (padding) \\ \hline
\end{tabular}
\end{table}

\subsection{Análisis de Modos de Operación}

La comparación visual entre ECB y CBC demostró claramente las diferencias de seguridad:

\begin{itemize}
    \item \textbf{ECB:} Los patrones repetitivos en el texto plano se reflejaban en el texto cifrado (estructuras visibles en hexdump)
    \item \textbf{CBC:} Los patrones no eran visibles gracias al encadenamiento de bloques
\end{itemize}

\subsection{Verificación de Integridad}

Los hashes MD5 del archivo original coincidieron exactamente con los del archivo descifrado, verificando que no hubo pérdida de datos durante el proceso de cifrado y descifrado.

% =========================================================
%                     ANÁLISIS
% =========================================================
\section{Análisis}

\subsection{Seguridad de AES}

AES-256 proporciona un nivel de seguridad muy alto contra ataques de fuerza bruta. Con $2^{256}$ posibles claves, incluso con computadoras cuánticamente avanzadas, el tiempo requerido para un ataque exitoso sería prohibitivamente largo.

\subsection{Importancia del Vector de Inicialización (IV)}

El IV aleatorio utilizado en CBC es crítico para la seguridad. Sin un IV único para cada cifrado, patrones en el texto plano podrían ser detectados, especialmente cuando se cifran múltiples archivos con la misma clave.

\subsection{Padding y Tamaño de Bloque}

Los archivos cifrados son ligeramente más grandes debido al padding PKCS\#7, que asegura que el texto plano sea un múltiplo del tamaño de bloque (128 bits para AES).

\subsection{Comparación ECB vs CBC}

ECB demoró ser inseguro para datos con patrones, mientras que CBC, al ser dependiente del bloque anterior, mantiene la aleatoriedad incluso con patrones de entrada. Esto confirma la necesidad de utilizar modos seguros en aplicaciones prácticas.

\subsection{Rendimiento}

OpenSSL ejecutó las operaciones de cifrado y descifrado de manera eficiente, incluso para archivos de tamaño considerable. No se observaron cuellos de botella significativos.

% =========================================================
%                  CONCLUSIONES
% =========================================================
\section{Conclusiones}

\begin{enumerate}
    \item OpenSSL proporciona herramientas robustas y eficientes para la implementación de criptografía simétrica en sistemas operativos Unix/Linux.

    \item AES-256 en modo CBC ofrece un nivel de seguridad prácticamente inquebrantable para datos sensibles, siendo la opción recomendada para protección moderna de datos.

    \item La generación correcta de claves aleatorias y vectores de inicialización es fundamental para mantener la seguridad del sistema criptográfico.

    \item El modo ECB, aunque simple, debe evitarse en producciones reales debido a la vulnerabilidad ante análisis de patrones. Modos como CBC, CTR u otros deben preferirse.

    \item La implementación práctica de criptografía requiere no solo conocimiento teórico, sino también familiaridad con herramientas específicas y buenas prácticas de seguridad.

    \item La verificación de integridad de datos mediante hash es una práctica complementaria esencial para detectar modificaciones no autorizadas.

    \item El cifrado de archivos debe ser parte de una estrategia más amplia de seguridad que incluya gestión de acceso, auditoría y monitoreo continuo.
\end{enumerate}

% =========================================================
%                   RECOMENDACIONES
% =========================================================
\section{Recomendaciones}

\begin{itemize}
    \item \textbf{Utilizar AES-256:} Preferir AES-256-CBC o AES-256-GCM en modo de autenticación para máxima seguridad.

    \item \textbf{Gestión segura de claves:} Almacenar claves en sistemas de gestión de claves (Key Management Services) como HashiCorp Vault o AWS KMS.

    \item \textbf{IV aleatorios:} Generar un IV nuevo para cada operación de cifrado, especialmente en modo CBC.

    \item \textbf{Autenticación:} Además de confidencialidad, implementar integridad mediante HMAC o AEAD (Authenticated Encryption with Associated Data).

    \item \textbf{Rotación de claves:} Implementar políticas de rotación periódica de claves (anualmente como mínimo).

    \item \textbf{Auditoría:} Mantener registros de acceso y uso de operaciones criptográficas para propósitos de auditoría y cumplimiento.
\end{itemize}

% =========================================================
%                   REFERENCIAS
% =========================================================
\newpage
\section*{Referencias}

\begin{enumerate}
    \item OpenSSL Documentation. (2024). \textit{OpenSSL Command Line Utility Documentation}.
    Disponible en: \url{https://www.openssl.org/docs/}

    \item NIST. (2001). \textit{Specification for the Advanced Encryption Standard (AES)}.
    FIPS 197.

    \item Stallings, W. (2017). \textit{Cryptography and Network Security: Principles and Practice} (7ª ed.).
    Pearson Education.

    \item Ferguson, N., Schneier, B., \& Kohno, T. (2010). \textit{Cryptography Engineering:
    Design Principles and Practical Applications}. Wiley Publishing.

    \item Menezes, A. J., van Oorschot, P. C., \& Vanstone, S. A. (1996).
    \textit{Handbook of Applied Cryptography}. CRC Press.

    \item Katz, J., \& Lindell, Y. (2014). \textit{Introduction to Modern Cryptography} (2ª ed.).
    Chapman and Hall/CRC.

    \item ISO/IEC 18033-3:2010. \textit{Encryption algorithms -- Part 3: Block ciphers}.

    \item OWASP. (2024). \textit{Cryptographic Storage Cheat Sheet}.
    Disponible en: \url{https://cheatsheetseries.owasp.org/}
\end{enumerate}

\end{document}
